% !TEX program = xelatex
\documentclass[5pt,a4paper]{article}

% --- 中文支持 ---
\usepackage[UTF8]{ctex}
\setCJKmainfont{SimSun}[BoldFont=SimHei,ItalicFont=KaiTi]
\setCJKsansfont{SimHei}

% --- 极限边距 ---
\usepackage[top=5mm,bottom=5mm,left=5mm,right=5mm]{geometry}
\pagestyle{empty}
\setlength{\parindent}{0pt}
\setlength{\parskip}{0.5pt}
\renewcommand{\baselinestretch}{0.88}

% --- 字体 ---
\usepackage{anyfontsize}
\usepackage{multicol}
\usepackage{pdfpages}
\setlength{\columnsep}{3.5mm}
\setlength{\columnseprule}{0.2pt}

% --- 颜色 ---
\usepackage[dvipsnames]{xcolor}
\definecolor{h1}{HTML}{0D47A1}
\definecolor{h2}{HTML}{1565C0}
\definecolor{h3}{HTML}{E65100}
\definecolor{warn}{HTML}{C62828}
\definecolor{tip}{HTML}{2E7D32}
\definecolor{bg}{HTML}{EEF2FF}
\definecolor{bgwarn}{HTML}{FFF3E0}
\definecolor{bgex}{HTML}{F0FFF0}

% --- 数学 ---
\usepackage{amsmath,amssymb}

% --- 紧凑列表 ---
\usepackage{enumitem}
\setlist{nosep,leftmargin=8pt,labelsep=2pt,itemsep=0pt,parsep=0pt,topsep=0pt}

% --- 自定义命令 ---
\newcommand{\chap}[1]{\vspace{0.8pt}{\fontsize{6.2}{6.8}\selectfont\bfseries\color{white}\colorbox{h1}{\strut\hspace{1pt}#1\hspace{1pt}}}\par\vspace{0.2pt}}
\newcommand{\sect}[1]{\vspace{0.2pt}{\fontsize{5.3}{5.8}\selectfont\bfseries\color{h2}#1}\par\vspace{0.05pt}}
\newcommand{\exm}[1]{\vspace{0.15pt}{\fontsize{5.2}{6.0}\selectfont\bfseries\color{tip}#1}\par\vspace{0.1pt}}
\newcommand{\key}[1]{{\bfseries\color{h3}#1}}
\newcommand{\warn}[1]{{\color{warn}\bfseries #1}}
\newcommand{\fml}[1]{\;\colorbox{bg}{$\displaystyle #1$}\;}
\newcommand{\wfml}[1]{\;\colorbox{bgwarn}{$\displaystyle #1$}\;}

\begin{document}
\includepdf[pages={1,2},fitpaper=true]{../../期中速查表.pdf}
\fontsize{6.2}{7.8}\selectfont

\iffalse
\begin{multicols}{3}

% ============================================================
\chap{第1章\;流体物理性质}
% ============================================================

\sect{连续介质假设}
流体质点：宏观$\infty$小、微观$\infty$大。
Kn$=\lambda/L\ll1$（$\lambda$=分子自由程，$L$=特征尺度）。
使物理量成坐标和时间的\key{连续函数}，可用微积分。

\sect{四大力学性质}
\key{易流动性}、\key{粘性}、\key{压缩性}、\key{表面张力}

\sect{粘性（核心）}
\key{牛顿内摩擦定律}：\fml{\tau = \mu \dfrac{\partial u}{\partial y}}

$\mu$=动力粘度(Pa$\cdot$s)；$\nu=\mu/\rho$=运动粘度(m$^2$/s)

\warn{$\mu$与$\nu$别搞混}！题目给$\nu$时先算$\mu=\rho\nu$

温度$\uparrow$：\key{液体}$\mu\downarrow$（引力$\downarrow$）；\key{气体}$\mu\uparrow$（热运动$\uparrow$）

牛顿流体：$\tau$-$\partial u/\partial y$线性，$\mu$与速度梯度\key{无关}

\sect{作用力}
\key{质量力}：$\propto m$，重力/惯性力/离心力\\
\key{表面力}：$\propto A$。法向=压力；切向=切应力

\sect{切应力计算}
(1)$du/dy$ (2)$\tau=\mu\cdot du/dy$ (3)壁面($y\!=\!0$)$\tau$最大 (4)代值

\exm{例：$u\!=\!C(50y\!-\!y^2)$，$\mu\!=\!0.05$}
$du/dy=C(50-2y)$，$\tau=0.05C(50-2y)$\\
$y\!=\!0$时$\tau_{\max}\!=\!0.05\!\times\!50C\!=\!2.5C$ Pa

\sect{表面张力}
水滴$\Delta p\!=\!2\sigma/R$；气泡$\Delta p\!=\!4\sigma/R$\\
润湿$\theta\!<\!90°$凹$\uparrow$；不润湿$\theta\!>\!90°$凸$\downarrow$

单位：1泊=0.1Pa$\cdot$s；1斯=$10^{-4}$m$^2$/s

\exm{例：滑动轴承功率}
$D$=12cm，$b$=20cm，间隙$t$=0.1cm，$\mu$已知，转速$n$\\
线速度$v=\pi Dn/60$，$\tau=\mu v/t$\\
摩擦力$F=\tau\!\cdot\!\pi Db$，功率$P=Fv$

\vspace{0.8pt}\hrule height 0.4pt\vspace{0.8pt}

% ============================================================
\chap{第2章\;流体静力学}
% ============================================================

\sect{基本方程}
\fml{p = p_0 + \rho g h}

\key{等压面}：连通静止同种流体，同深度$p$相等\\
\warn{不连通或不同液体时不成立}

\sect{三种压强}
$p_{abs}\!=\!p_a+p_{gage}$；$p_{vac}\!=\!p_a-p_{abs}\!=\!-p_{gage}$\\
\warn{大题用表压}！开口水面表压$p_0\!=\!0$

\sect{U管测压——爬楼梯法}
从已知点出发，\key{下$+\rho gh$，上$-\rho gh$}，遇不同液体用等压面跨过

\exm{例：复式U管测压}
从A点出发逐段推：$p_A+\rho_1 g h_1-\rho_2 g h_2+\rho_3 g h_3-\cdots=p_B$

\sect{平面总压力}
合力：\fml{F = \rho g h_C \cdot A}

$h_C$=形心\key{垂直}深度，$A$=面积

压力中心：\fml{y_D = y_C + \frac{J_C}{y_C A}}

$y_C$=形心沿斜面距离；$J_C$=对形心轴惯性矩\\
矩形$J_C\!=\!bh^3\!/12$；圆$J_C\!=\!\pi d^4\!/64$；三角形$J_C\!=\!bh^3\!/36$\\
\key{压力中心永远在形心下方}

\exm{例：矩形平板受力}
宽$b$=2m，高$h$=3m，顶边淹深$H$=1m\\
形心深度$h_C\!=\!H+h/2\!=\!2.5$m\\
$F\!=\!\rho g h_C\!\cdot\!bh\!=\!1000\!\times\!9.81\!\times\!2.5\!\times\!6\!=\!147.15$kN\\
$y_C\!=\!h_C\!=\!2.5$m（垂直板），$J_C\!=\!2\!\times\!3^3\!/12\!=\!4.5$m$^4$\\
$y_D\!=\!2.5+4.5/(2.5\!\times\!6)\!=\!2.8$m

\sect{曲面总压力（压轴大题）}
\key{水平分力}：\fml{F_x = \rho g h_{Cx} \cdot A_x}

$A_x$=曲面在\key{铅直面投影面积}，$h_{Cx}$=投影形心深度

\key{竖直分力}：\fml{F_y = \rho g V_{\text{压}}}

压力体：曲面边界\key{竖直向上}画到液面围成的体积\\
\key{实}压力体（水在上压）→$F_y$\key{向下}\\
\key{虚}压力体（水在下托）→$F_y$\key{向上}

合成：$F\!=\!\sqrt{F_x^2\!+\!F_y^2}$，$\tan\theta\!=\!F_y/F_x$

\warn{圆柱面：合力作用线必过圆心！}

\exm{例：圆弧闸门$r$=0.8m，宽$L$=1m}
轴在液面下$h$=2.4m，圆心角$90°$\\
$A_x\!=\!r\!\times\!L\!=\!0.8$m$^2$，$h_{Cx}\!=\!h+r/2\!=\!2.8$m\\
$F_x\!=\!1000\!\times\!9.81\!\times\!2.8\!\times\!0.8\!=\!21.97$kN\\
$V_\text{压}\!=\!h\!\cdot\!r\!\cdot\!L-\frac{1}{4}\pi r^2 L\!=\!(1.92\!-\!0.503)\!=\!1.417$m$^3$（或含加减）\\
$F_y\!=\!\rho g V_\text{压}$，合力$F\!=\!\sqrt{F_x^2+F_y^2}$\\
作用线过圆心，$\theta\!=\!\arctan(F_y/F_x)$

\sect{易错点}
\warn{(1)}忘乘宽度$L$！ \warn{(2)}封闭容器等效液面升高$\Delta h\!=\!p_0/(\rho g)$ \warn{(3)}$h_C$=垂直深度$\neq$斜距

\sect{相对平衡}
匀加速：等压面倾斜$\tan\alpha\!=\!a_x/(g+a_z)$\\
等角速旋转：$z\!=\!\omega^2r^2/(2g)+C$（旋转抛物面）\\
\key{旋转抛物体体积=同高圆柱$\times\frac{1}{2}$}

\exm{例：圆柱容器旋转}
高1.8m，直径0.9m，水深1.35m\\
无溢出最大转速：液面顶端到容器边缘\\
$\Delta z=\omega^2 R^2/(2g)=2(H-h_0)$→解出$\omega$→$n=60\omega/(2\pi)$\\
利用旋转抛物体体积=柱体$\times1/2$

\vspace{0.8pt}\hrule height 0.4pt\vspace{0.8pt}

% ============================================================
\chap{第3章\;流体运动学}
% ============================================================

\sect{两种方法}
\key{拉格朗日}：盯质点追踪→迹线 | \key{欧拉}：固定空间观察（常用）\\
PIV→\key{欧拉法}

\sect{定常 vs 非定常}
定常：$\partial\vec{v}/\partial t\!=\!0$ | \warn{定常$\neq$无加速度}（迁移加速度$\neq 0$）\\
等速直线$\neq$定常 | 定常中质点速度未必与$t$无关

\sect{流线·迹线·脉线}
流线=某瞬间与$\vec{v}$相切（\key{不相交}）\\
迹线=质点轨迹 | 脉线=经同一点的质点连线\\
\warn{定常流三线合一}；一维流也三线合一

\sect{物质导数}
\fml{\frac{D}{Dt}=\frac{\partial}{\partial t}+u\frac{\partial}{\partial x}+v\frac{\partial}{\partial y}+w\frac{\partial}{\partial z}}
= 当地导数 + 迁移导数

\sect{加速度公式}
\fml{a_x=u_t+u u_x+v u_y+w u_z}
\fml{a_y=v_t+u v_x+v v_y+w v_z}

\warn{$a_x$中有$v\cdot\partial u/\partial y$，$a_y$中有$u\cdot\partial v/\partial x$，别搞反！}

\sect{加速度四步法}
(1)辨识$u,v$ (2)求6个偏导 (3)组装 (4)代入$(x_0,y_0,t_0)$

\exm{例1：$u\!=\!2x\!+\!t$，$v\!=\!-2y\!+\!3$，$t\!=\!1$，点$(1,2)$}
$\partial u/\partial t\!=\!1$，$\partial u/\partial x\!=\!2$，$\partial u/\partial y\!=\!0$\\
$\partial v/\partial t\!=\!0$，$\partial v/\partial x\!=\!0$，$\partial v/\partial y\!=\!-2$\\
$a_x\!=\!1+(2x\!+\!t)\!\cdot\!2+(-2y\!+\!3)\!\cdot\!0\!=\!1+2u$\\
$a_y\!=\!0+(2x\!+\!t)\!\cdot\!0+(-2y\!+\!3)\!\cdot\!(-2)\!=\!-2v$\\
$t\!=\!1$: $u(1,2,1)\!=\!3$, $v(1,2,1)\!=\!-1$\\
$a_x\!=\!1+6\!=\!7$，$a_y\!=\!-2\!\times\!(-1)\!=\!2$\\
$a=\sqrt{49+4}=\sqrt{53}\approx7.28$ m/s$^2$

\exm{例2：$u\!=\!3x^2$，$v\!=\!-6xy$，点$(2,1)$}
$\partial u/\partial x\!=\!6x$, $\partial u/\partial y\!=\!0$, $\partial v/\partial x\!=\!-6y$, $\partial v/\partial y\!=\!-6x$\\
$a_x\!=\!3x^2\!\cdot\!6x\!=\!18x^3\!=\!144$\\
$a_y\!=\!3x^2\!\cdot\!(-6y)+(-6xy)\!\cdot\!(-6x)\!=\!18x^2y\!=\!72$\\
$a\!=\!\sqrt{144^2\!+\!72^2}\!\approx\!161.0$ m/s$^2$

\sect{流线方程}
\fml{\frac{dx}{u}=\frac{dy}{v}}
交叉相乘→分离变量→积分→流线族

\warn{非定常：先代$t\!=\!t_0$再积分}；保留$C$

\exm{例：$u\!=\!2x$，$v\!=\!-2y$}
$dx/(2x)\!=\!dy/(-2y)$→$dx/x\!=\!-dy/y$\\
$\ln|x|\!=\!-\ln|y|+C_1$→\fml{xy=C}（等轴双曲线）

\exm{例：$u\!=\!Ay$，$v\!=\!Ax$}
$dx/(Ay)\!=\!dy/(Ax)$→$xdx\!=\!ydy$\\
$x^2/2\!=\!y^2/2+C_1$→\fml{x^2-y^2=C}

\exm{齐次方程法：$u\!=\!x\!+\!y$，$v\!=\!x\!-\!y$}
$(x\!-\!y)dx\!=\!(x\!+\!y)dy$，令$y\!=\!sx$→\fml{x^2-2xy-y^2=C}

\sect{微团运动分解}
\key{4部分}：平移、旋转、线变形、角变形

$\omega_z\!=\!\frac{1}{2}(\frac{\partial v}{\partial x}-\frac{\partial u}{\partial y})$；
涡量$\vec{\Omega}\!=\!\nabla\times\vec{v}\!=\!2\vec{\omega}$

无旋$\iff\nabla\times\vec{v}\!=\!0$；\warn{直线$\neq$无旋}

线变形率$\varepsilon_{xx}\!=\!\partial u/\partial x$，$\varepsilon_{yy}\!=\!\partial v/\partial y$
角变形率$\varepsilon_{xy}\!=\!\frac{1}{2}(\partial u/\partial y+\partial v/\partial x)$

\exm{例：$u\!=\!xy$，$v\!=\!x^2\!-\!y^2/2$，在$(2,2)$}
$\omega_z\!=\!\frac{1}{2}(2x-x)\!=\!x/2\!=\!1$，有旋\\
$\varepsilon_{xx}\!=\!y\!=\!2$，$\varepsilon_{yy}\!=\!-y\!=\!-2$（不可压$\checkmark$）\\
$\varepsilon_{xy}\!=\!\frac{1}{2}(x+2x)\!=\!3x/2\!=\!3$

\sect{迹线方程求法}
$\frac{dx}{dt}\!=\!u$，$\frac{dy}{dt}\!=\!v$→分别积分

\exm{例：$u\!=\!t$，$v\!=\!2t$，初始$(0,0)$}
$dx/dt\!=\!t$→$x\!=\!t^2/2+C_1$，$x(0)\!=\!0$→$C_1\!=\!0$\\
$dy/dt\!=\!2t$→$y\!=\!t^2+C_2$，$y(0)\!=\!0$→$C_2\!=\!0$\\
消参：$y\!=\!2x$（直线）

\exm{例：$u\!=\!2x$，$v\!=\!2y$，$t_0\!=\!0$}
$dx/dt\!=\!2x$→$x\!=\!x_0 e^{2t}$\\
$dy/dt\!=\!2y$→$y\!=\!y_0 e^{2t}$\\
消参$t$：$x/x_0\!=\!y/y_0$→$y\!=\!(y_0/x_0)x$（直线族）\\
此为定常流→与流线$y/x\!=\!C$完全一致$\checkmark$

\vspace{0.8pt}\hrule height 0.4pt\vspace{0.8pt}

% ============================================================
\chap{Ch1-3综合速查}
% ============================================================

\sect{常用面积/体积/惯性矩}
矩形：$A\!=\!bh$，$J_C\!=\!bh^3/12$\\
圆：$A\!=\!\pi d^2/4$，$J_C\!=\!\pi d^4/64$\\
三角形：$A\!=\!bh/2$，$J_C\!=\!bh^3/36$（对底边平行形心轴）\\
半圆：$A\!=\!\pi R^2/2$，形心距底$\bar{y}\!=\!4R/(3\pi)$

圆柱：$V\!=\!\pi R^2 h$ | 球：$V\!=\!4\pi R^3/3$ | 锥：$V\!=\!\pi R^2h/3$\\
圆柱面投影面积=直径$\times$高=$dh$

\sect{判断题速查}
$\bullet$ 静止流体中只存在\key{压力}（无切应力）→$\checkmark$\\
$\bullet$ 理想流体=无粘性，$\neq$不可压缩\\
$\bullet$ 粘性系数只随\key{温度}变化→$\checkmark$\\
$\bullet$ 均质不可压流体静止$\to$质量力有势→$\checkmark$\\
$\bullet$ 等压面上各点压强相等必须\key{同一种连续静止}流体\\
$\bullet$ ``理想流体定常流的连续性方程为$\nabla\!\cdot\!\vec{v}\!=\!0$''→\warn{$\times$}\\
$\bullet$ ``系统的表面会不断变形''→$\checkmark$\\
$\bullet$ ``定常流动中流体的加速度为零''→\warn{$\times$}有迁移加速度\\
$\bullet$ ``一维流动中流线、迹线、脉线重合''→$\checkmark$\\
$\bullet$ 直线运动的流体$\neq$一定无旋

\sect{典型公式快查}
$\rho_{\text{水}}\!=\!1000$ kg/m$^3$ | $g\!=\!9.81$ m/s$^2$（或10）\\
$1$atm$\!=\!101325$Pa$\!=\!10.33$mH$_2$O$\!=\!760$mmHg

圆管$A\!=\!\pi d^2/4$；$d$小一半→$A$变$1/4$→$V$变$4$倍

伯努利压强形式：$p+\frac{1}{2}\rho v^2+\rho gz\!=\!$const

动压$q\!=\!\frac{1}{2}\rho v^2$（伯努利压强差项）

\sect{平衡微分方程（欧拉静力学）}
$\frac{\partial p}{\partial x}\!=\!\rho f_x$，$\frac{\partial p}{\partial y}\!=\!\rho f_y$，$\frac{\partial p}{\partial z}\!=\!\rho(f_z)$

重力场中$f_x\!=\!f_y\!=\!0$，$f_z\!=\!-g$→$dp\!=\!-\rho g\,dz$

\sect{考试计算题结构}
\key{第1题}(约10分)：求平面/曲面总压力\\
\key{第2题}(约8分)：加速度或流线计算\\
\key{第3题}(约10分)：伯努利+连续方程联合求解\\
\key{第4题}(约8分)：动量方程求弯管/平板受力\\
填空/判断约30分，选择约10分

\sect{压力体体积速算}
$\frac{1}{4}$圆柱=$\frac{1}{4}\pi R^2 L$ | $\frac{1}{2}$圆柱=$\frac{1}{2}\pi R^2 L$\\
长方体$-\frac{1}{4}$圆柱=$rL(h-\frac{\pi r}{4})$（常见压力体）

圆弧闸门典型：上方长方体$V_1\!=\!rLh$减去$\frac{1}{4}$圆柱$V_2\!=\!\frac{\pi r^2L}{4}$

\sect{计算注意事项}
$\bullet$ 所有面积计算别忘$\pi$和宽度$L$\\
$\bullet$ $\sin30°\!=\!0.5$, $\cos30°\!=\!\frac{\sqrt{3}}{2}$, $\tan45°\!=\!1$\\
$\bullet$ $\sin60°\!=\!\frac{\sqrt{3}}{2}$, $\cos60°\!=\!0.5$, $\tan60°\!=\!\sqrt{3}$\\
$\bullet$ $\frac{\partial}{\partial x}(x^n)\!=\!nx^{n-1}$ | $\int x^n dx\!=\!\frac{x^{n+1}}{n+1}$\\
$\bullet$ $\frac{\partial}{\partial x}(e^{ax})\!=\!ae^{ax}$ | $\frac{\partial}{\partial x}(\ln x)\!=\!1/x$\\
$\bullet$ $\int\frac{1}{x}dx\!=\!\ln|x|+C$ | $\int\frac{1}{x^2}dx\!=\!-\frac{1}{x}+C$\\
$\bullet$ $\int 2\pi r\,dr\!=\!\pi r^2$ | $\int_0^R r(1\!-\!r^2/R^2)dr\!=\!R^2/4$

\sect{倾斜平面计算要点}
$h_C$=形心\key{垂直}深度；$y_C$=形心沿\key{斜面}距离\\
两者关系：$h_C\!=\!y_C\sin\theta$（$\theta$=斜面与水平夹角）\\
正对液面的垂直板：$h_C\!=\!y_C$，$\theta\!=\!90°$

\sect{多层液体问题}
不同密度液体分层→分段算压强\\
界面处压强连续：$p_0+\rho_1 g h_1\!=\!$界面$p$\\
再往下：$p_{\text{界面}}+\rho_2 g h_2$

\sect{容器匀加速运动}
水平加速$a_x$：等压面倾角$\tan\alpha\!=\!a_x/g$\\
竖直加减速：$p\!=\!\rho(g\pm a)h$——向上加速用$g\!+\!a$，向下用$g\!-\!a$\\
自由落体$a\!=\!g$→$p\!=\!0$（失重）

\exm{例：水平加速度$a$=3.27m/s$^2$}
$\tan\alpha\!=\!3.27/9.81\!=\!1/3$\\
液面倾斜角$\alpha\!=\!\arctan(1/3)\!=\!18.4°$\\
前壁增深$\Delta h\!=\!L\tan\alpha/2$，后壁减深同量

\sect{斯托克斯定理}
$\oint_C \vec{v}\!\cdot\!d\vec{l}\!=\!\iint_S (\nabla\!\times\!\vec{v})\!\cdot\!d\vec{A}$\\
速度环量$\Gamma\!=\!\oint\vec{v}\!\cdot\!d\vec{l}$\\
开尔文定理：正压理想流体$\Gamma$不随$t$变

\sect{流量计算集锦}
矩形截面$A\!=\!bh$，$Q\!=\!Abv$\\
圆截面$A\!=\!\pi d^2/4$，$Q\!=\!\pi d^2v/4$\\
梯形截面$A\!=\!(b_1\!+\!b_2)h/2$\\
层流圆管$Q\!=\!\pi R^4\Delta p/(8\mu L)$（哈根-泊肃叶）

\end{multicols}

\clearpage
\begin{multicols}{3}

% ============================================================
\chap{第4章\;RTT与连续方程}
% ============================================================

\sect{系统 vs 控制体}
\key{系统}：同一批质点，质量不变，边界会变形\\
\key{控制体}(CV)：固定空间区域，流体穿越边界进出\\
控制面(CS)：进流面/出流面/固壁面

\sect{雷诺输运方程}
\fml{\frac{DB_\text{sys}}{Dt}=\frac{\partial}{\partial t}\!\int_{CV}\!\rho b\,dV+\oint_{CS}\!\rho b(\vec{v}\!\cdot\!\vec{n})\,dA}

左=系统随体变化率；右=\key{储存变化+净流出}\\
$\vec{n}$外法向：流出$+$，流入$-$

\key{必背}：系统随体导数 = CV内偏$t$ + CS净通量

\sect{连续性方程（$B\!=\!m$, $b\!=\!1$）}
积分：$\frac{\partial}{\partial t}\!\int\!\rho dV+\oint\!\rho(\vec{v}\!\cdot\!\vec{n})dA\!=\!0$

微分：\fml{\frac{\partial\rho}{\partial t}+\nabla\!\cdot\!(\rho\vec{v})=0}

\sect{简化形式}
不可压缩：\wfml{\nabla\!\cdot\!\vec{v}=0} 即$\frac{\partial u}{\partial x}+\frac{\partial v}{\partial y}+\frac{\partial w}{\partial z}=0$

定常可压：$\nabla\!\cdot\!(\rho\vec{v})\!=\!0$ | 一维定常不可压：\fml{A_1V_1\!=\!A_2V_2\!=\!Q}

\warn{$\nabla\!\cdot\!\vec{v}\!=\!0$表示不可压，$\neq$理想流体定常！}

\sect{连续性验证}
验$\partial u/\partial x+\partial v/\partial y=0$?

\exm{例：$u\!=\!x^2\!+\!y$，$v\!=\!-2xy\!+\!3t$}
$\partial u/\partial x\!=\!2x$，$\partial v/\partial y\!=\!-2x$→$2x\!-\!2x\!=\!0\;\checkmark$

\exm{由连续方程求$v$：已知$u\!=\!3x^2y$}
$\partial u/\partial x\!=\!6xy$→$\partial v/\partial y\!=\!-6xy$\\
$v\!=\!\int(-6xy)dy\!=\!-3xy^2+f(x)$

\sect{管道变截面}
$V_2\!=\!V_1 d_1^2/d_2^2$；直径缩半→速度$\times4$

\exm{例：$D$=200mm→$d$=100mm，$V_1$=2m/s}
$V_2\!=\!2\!\times\!200^2/100^2\!=\!8$ m/s

\sect{水箱液面变化率}
\fml{\frac{dh}{dt}=\frac{A_1V_1-A_2V_2}{A_0}}

封闭有气：$\frac{dh}{dt}\!=\!\frac{\rho(A_1V_1\!-\!A_2V_2)}{(\rho\!-\!\rho_a)A_0}$

\exm{例：$D_0$=2m，出水管$d$=0.1m，$V_{out}$=4m/s，$Q_{in}$=0.02m$^3$/s}
$Q_{out}\!=\!\pi\!\times\!0.01/4\!\times\!4\!=\!0.0314$\\
$dh/dt\!=\!(0.02\!-\!0.0314)/\pi\!\approx\!-3.6$mm/s（$\downarrow$）

\sect{抛物线速度求流量}
$v\!=\!v_0(1\!-\!r^2\!/R^2)$→$Q\!=\!\frac{\pi R^2 v_0}{2}$

\warn{非均匀分布必须积分！}层流$\bar{v}\!=\!v_0/2$

\sect{散度物理含义}
$\nabla\!\cdot\!\vec{v}$=单位体积膨胀率\\
$\nabla\!\cdot\!\vec{v}\!>\!0$：流体膨胀（源）；$<\!0$：收缩（汇）\\
不可压$\nabla\!\cdot\!\vec{v}\!=\!0$：无净体积变化

\sect{RTT三大方程对照}
$B\!=\!m,b\!=\!1$→\key{连续性方程}\\
$B\!=\!m\vec{v},b\!=\!\vec{v}$→\key{动量方程}\\
$B\!=\!E,b\!=\!e$→\key{能量方程}

\sect{动量方程（积分形式）}
\fml{\sum\vec{F}=\frac{\partial}{\partial t}\!\int_{CV}\rho\vec{v}\,dV+\oint_{CS}\rho\vec{v}(\vec{v}\!\cdot\!\vec{n})\,dA}

定常一进一出常写成：$\sum\vec{F}\!=\!\dot{m}(\beta_2\vec{v}_2-\beta_1\vec{v}_1)$\\
\warn{外力=压强力+质量力+壁面反力；速度/力都要按\key{矢量}列分量}

\sect{多管汇流/分流}
$Q_{\text{总}}\!=\!\sum Q_i$；$A_1V_1\!=\!A_2V_2+A_3V_3$

\sect{不连续判断举例}
$u\!=\!x+y^2,v\!=\!x^2\!-\!y$: $1-1\!=\!0\;\checkmark$\\
$u\!=\!2xy,v\!=\!x^2\!+\!y^2$: $2y+2y\!=\!4y\;\times$

\vspace{0.8pt}\hrule height 0.4pt\vspace{0.8pt}

% ============================================================
\chap{第4章补充\;应力张量/微分方程}
% ============================================================

\sect{任意方位面的应力}
\fml{\vec{p}_n=\vec{n}\cdot\mathbf{P}}

$\vec{p}_n$通常\key{不与}$\vec{n}$共线；法向应力$p_{nn}\!=\!\vec{p}_n\!\cdot\!\vec{n}$\\
切向大小：$p_{n\tau}\!=\!\sqrt{|\vec{p}_n|^2-p_{nn}^2}$

\sect{二阶应力张量}
\fml{\mathbf{P}\!=\!\left[\begin{smallmatrix}p_{xx}&p_{xy}&p_{xz}\\ p_{yx}&p_{yy}&p_{yz}\\ p_{zx}&p_{zy}&p_{zz}\end{smallmatrix}\right]}

动量矩守恒→\key{对称}：$p_{xy}\!=\!p_{yx}$，$p_{yz}\!=\!p_{zy}$，$p_{zx}\!=\!p_{xz}$\\
故9个分量仅6个独立；$\nabla\!\cdot\!\mathbf{P}$=单位体积表面力

\sect{表面力体积化}
\fml{\iint_{CS}\vec{p}_n\,dA=\iiint_{CV}\nabla\!\cdot\!\mathbf{P}\,dV}

微分法研究的是\key{全场分布}$(u,v,w,p)$，不是截面平均量\\
\warn{题目一旦问“空间分布/偏微分方程组/定解条件”就该想到这一套}

\sect{静止/理想流体应力状态}
静止流体：\fml{\mathbf{P}=-p\mathbf{I}}

故任意面应力$\vec{p}_n\!=\!-p\vec{n}$，\key{只有压应力、无切应力}\\
无粘理想流体同样无切应力；粘性流体切应力来自速度梯度

\sect{动量微分型基本方程}
\fml{\rho\frac{D\vec{v}}{Dt}=\rho\vec{f}+\nabla\!\cdot\!\mathbf{P}}

再配合连续方程$\frac{\partial\rho}{\partial t}+\nabla\!\cdot(\rho\vec{v})\!=\!0$才能闭合\\
未知常为$u,v,w,p$；若可压还要配状态方程/能量方程

\sect{欧拉方程}
理想流体$\mathbf{P}\!=\!-p\mathbf{I}$→\fml{\rho\frac{D\vec{v}}{Dt}=-\nabla p+\rho\vec{f}}

若$\vec{f}\!=\!\vec{g}$，沿流线可写：\fml{\frac{dp}{\rho}+V\,dV+g\,dz=0}\\
对定常无粘不可压积分后，才得到伯努利方程

\sect{Navier-Stokes}
牛顿流体：$\tau_{ij}\!=\!\mu\!\left(\partial u_i/\partial x_j+\partial u_j/\partial x_i\right)$\\
不可压$\mu$常数：\fml{\rho\frac{D\vec{v}}{Dt}=-\nabla p+\mu\nabla^2\vec{v}+\rho\vec{f}}

\warn{无粘极限→欧拉；低Re粘性项常不可忽略}

\sect{边界与定解条件}
初始条件：$t\!=\!0$时速度/压强分布\\
粘性固壁：\key{无滑移}$\vec{v}_L\!=\!\vec{v}_{wall}$；无粘固壁：只需法向速度相等\\
入口常给速度，出口常给压强/充分发展条件

\vspace{0.8pt}\hrule height 0.4pt\vspace{0.8pt}

% ============================================================
\chap{第4章续\;伯努利/能量/动量}
% ============================================================

\sect{元流伯努利方程}
\fml{\frac{p_1}{\rho g}+\frac{v_1^2}{2g}+z_1=\frac{p_2}{\rho g}+\frac{v_2^2}{2g}+z_2}

\key{7条件}：理想/重力/定常/绝热/无轴功/不可压/同一流线

$z$=位置水头，$p/(\rho g)$=压力水头，$v^2/(2g)$=速度水头\\
总水头$H$=const沿流线；测压管水头$=z+p/(\rho g)$

\sect{理想流体能量方程}
\fml{h+\frac{V^2}{2}+gz=\text{const}}

适用于定常/无粘/质量力有势/绝热/沿流线\\
可压流优先记这条；不可压且$dh\!=\!dp/\rho$时才退化为伯努利

\sect{欧拉积分到伯努利}
\fml{\frac{dp}{\rho}+V\,dV+g\,dz=0}

积分得：$p/\rho+V^2/2+gz=\text{const}$；再除$g$得水头式\\
静止流体$V\!=\!0$时退化为$p/(\rho g)+z=\text{const}$

\sect{总流伯努利（含损失）}
\fml{z_1\!+\!\frac{p_1}{\rho g}\!+\!\frac{\alpha_1v_1^2}{2g}\!=\!z_2\!+\!\frac{p_2}{\rho g}\!+\!\frac{\alpha_2v_2^2}{2g}\!+\!h_w}

$\alpha$=动能修正：层流$\alpha\!=\!2$，湍流$\alpha\!\approx\!1.05$\\
$h_w\!=\!h_f+h_j$；含泵加$H_P$，含水轮机减$H_T$

\sect{总压/驻点/跨流线}
总压（驻点压）：\fml{p_0=p+\frac{1}{2}\rho v^2}

驻点$v\!=\!0$→该点压强就等于$p_0$；皮托管测的就是总压\\
\warn{有旋流只能沿同一流线用伯努利；无旋流才可跨流线}

\sect{常见特例}
大水箱液面：$p\!=\!0$，$v\!\approx\!0$；自由射流出口：$p\!=\!0$（表压）\\
水平管：$z_1\!=\!z_2$；等径管：$v_1\!=\!v_2$\\
速度增大$\Rightarrow$静压降低；总压在理想流中保持不变

\sect{解题步骤}
(1)选截面（一已知一求解）(2)选基准面\\
(3)列伯努利+连续方程 (4)$p$用\key{表压}，$v$用平均速度

\exm{例：收缩管$D$→$d$，$p_1$=$p_2$+$\Delta p$已知}
连续：$V_1A_1\!=\!V_2A_2$→$V_1\!=\!V_2 d^2/D^2$\\
伯努利（水平无损失）：$\frac{p_1}{\rho g}+\frac{V_1^2}{2g}=\frac{p_2}{\rho g}+\frac{V_2^2}{2g}$\\
代入解$V_2$→$Q\!=\!V_2\!\cdot\!\pi d^2/4$

\exm{例：水箱侧壁小孔出流}
取(1)液面(2)孔口出口。液面$V_1\!\approx\!0$，$p_1\!=\!0$\\
孔口$p_2\!=\!0$，$z_1\!-\!z_2\!=\!h$\\
$h\!=\!V_2^2/(2g)$→$V_2\!=\!\sqrt{2gh}$（托里拆利公式）

\sect{皮托管与文丘里管}
皮托管测\key{总压}$p_0\!=\!p+\frac{1}{2}\rho v^2$\\
静压管测\key{静压}$p$（壁面开孔）\\
流速：\fml{v=\sqrt{\frac{2(p_0-p)}{\rho}}=\sqrt{\frac{2\Delta p}{\rho}}}

U形管差压计：$\Delta p\!=\!(\rho_m\!-\!\rho)g\Delta h$\\
$v\!=\!\sqrt{2(\rho_m\!-\!\rho)g\Delta h/\rho}$

\key{文丘里管}：收缩段+喉部，利用$\Delta p$测流量\\
$Q\!=\!\frac{\pi\sqrt{2g}\,d_1^2 d_2^2}{4\sqrt{d_1^4-d_2^4}}\!\sqrt{h}\!=\!k\sqrt{h}$；实测$Q\!=\!\mu k\sqrt{h}$

\sect{孔口出流}
\fml{v_C\!=\!\varphi\sqrt{2gH_0};\quad Q\!=\!\mu A\sqrt{2gH_0}}
$\varphi$=速度系数$\approx$0.97；$\varepsilon$=收缩系数$\approx$0.64\\
$\mu\!=\!\varphi\varepsilon\!\approx\!0.62$（薄壁孔口）\\
$H_0\!=\!\alpha_0v_0^2/(2g)+H$（有来流速度时）\\
\warn{管嘴}：$\varepsilon\!=\!1,\mu\!\approx\!0.82$，流量系数比孔口\key{大}

\sect{动量方程}
\fml{\sum\vec{F}=\rho Q(\beta_2\vec{v}_2-\beta_1\vec{v}_1)}

$\beta$=动量修正：层流$\beta\!=\!4/3$，湍流$\beta\!\approx\!1.02$

\warn{$\vec{v}$带方向！先定正方向}

\exm{例：水柱冲击垂直平板}
来流$Q$，速度$V$，截后分两侧→板受力$F\!=\!\rho QV$

\exm{例：水柱被平板截取}
流量$Q$，来流$V$，截取$Q_1$，偏转角$\theta$\\
$x$：$F_x\!=\!\rho Q V-\rho(Q\!-\!Q_1)V\cos\theta$\\
$y$：$F_y\!=\!\rho(Q\!-\!Q_1)V\sin\theta$\\
联合解$F$和$\theta$

\exm{例：弯管受力}
选CV包围弯管段，进出截面1-2\\
$x$：$p_1A_1-p_2A_2\cos\alpha-F_{Rx}\!=\!\rho Q(V_2\cos\alpha-V_1)$\\
$y$：$-p_2A_2\sin\alpha+F_{Ry}\!=\!\rho Q V_2\sin\alpha$\\
解$F_R$即支座反力

\sect{动量矩}
$M\!=\!\rho Q(v_2r_2-v_1r_1)$

\sect{三种压强关系}
静压$p$+动压$\frac{1}{2}\rho v^2$=总压$p_0$\\
测压管水头$=z+p/(\rho g)$\\
总水头$H=z+p/(\rho g)+v^2/(2g)$=const

\sect{伯努利能量线图}
\key{EL}(总水头线)沿流程恒定/下降；\key{HGL}=EL$-v^2/(2g)$\\
截面缩小→$v\uparrow$→HGL$\downarrow$；\warn{HGL低于管轴}→负压

\exm{例：U管+伯努利（07真题型）}
油$\rho\!=\!920$，$d_A\!=\!100$mm，$d_B\!=\!200$mm，A静压管B皮托管\\
U管$\rho_m\!=\!13600$，$\Delta h\!=\!0.3$m。连续$v_A\!=\!4v_B$\\
伯努利：$p_A-p_B\!=\!\frac{15}{2}\rho v_B^2$\\
U管：$8\rho v_B^2\!=\!(\rho_m\!-\!\rho)g\Delta h$→$v_B\!\approx\!2.25$m/s

\vspace{0.8pt}\hrule height 0.4pt\vspace{0.8pt}

\end{multicols}

\clearpage
\fi

\begin{multicols}{3}

% ============================================================
\chap{第4章母题\;RTT/伯努利/动量}
% ============================================================

{\fontsize{5.4}{6.1}\selectfont

\sect{三年真题最常考}
06/07/08反复出现：\key{粘性拖板/轴承}、\key{U管测压}、\key{平面/曲面总压力}、\key{加速度/流线/有旋}、\key{旋转容器}、\key{RTT/液面变化率}、\key{伯努利+连续}、\key{皮托/泵}、\key{动量受力}。\\
提分顺序：\key{Ch2曲面压力}>\key{Ch5伯努利联立}>\key{Ch3加速度流线}>\key{Ch4连续}>\key{Ch1粘性}

\sect{母题1\;粘性拖板/滑动轴承}
两板距$h$、面积$A$、相对速差$U$：\fml{\tau=\mu U/h,\;F=\tau A,\;P=FU}

滑动轴承：$U=\pi Dn/60$，受剪面积$A=\pi Db$，故\\
\fml{P=\mu\frac{\pi Dn}{60t}\cdot \pi Db \cdot \frac{\pi Dn}{60}}
\warn{间隙很小时默认线性分布；题给$\nu$要先算$\mu=\rho\nu$}

\sect{母题2\;U管/封闭容器}
\key{爬楼梯法}：向下$+\rho gh$，向上$-\rho gh$，同一静止连通液体同水平面压强相等\\
$p_{abs}=p_a+p_g$，$p_{vac}=p_a-p_{abs}=-p_g$\\
封闭气腔等效液面：\fml{h_e=p_0/(\rho g)}，后面一律对\key{假想液面}算压力

\sect{母题3\;平面与曲面总压力}
平面：\fml{F=\rho g h_C A,\qquad y_D=y_C+\frac{I_C}{y_CA}}
\warn{$h_C$是\key{垂直深度}，$y_C$是沿斜面距离；压力中心总在形心下方}

曲面：\fml{F_x=\rho g h_{Cx}A_x,\qquad F_y=\rho gV_{\text{压}}}
实压力体$\Rightarrow F_y$向下；虚压力体$\Rightarrow F_y$向上；圆柱曲面\key{合力作用线过圆心}

\sect{母题4\;加速度/有旋/变形率}
\fml{a_x=u_t+uu_x+vu_y+wu_z,\qquad a_y=v_t+uv_x+vv_y+wv_z}
\fml{\omega_z=\frac{1}{2}(v_x-u_y),\quad \Omega_z=v_x-u_y}
\fml{\varepsilon_{xx}=u_x,\;\varepsilon_{yy}=v_y,\;\varepsilon_{xy}=\frac{u_y+v_x}{2}}
例：$u=2x+t,v=-2y+3$，在$(1,2,1)$：$u=3,v=-1$，故$a_x=7,a_y=2$

\sect{母题5\;流线/迹线}
流线：\fml{\frac{dx}{u}=\frac{dy}{v}=\frac{dz}{w}}；迹线：\fml{\frac{dx}{dt}=u,\;\frac{dy}{dt}=v}
非定常流求流线时，先把$t=t_0$代入速度场\\
例：$u=2x,v=-2y$，流线$\frac{dx}{2x}=\frac{dy}{-2y}\Rightarrow xy=C$；过$(2,2)$即$xy=4$

\sect{母题6\;连续方程/液面变化率}
\fml{\frac{\partial \rho}{\partial t}+\nabla\cdot(\rho\vec v)=0}
不可压：$u_x+v_y+w_z=0$；一维定常不可压：\fml{A_1V_1=A_2V_2=Q}
水箱液面变化：\fml{\frac{dh}{dt}=\frac{A_1V_1-A_2V_2}{A_0}}
抛物线速度$u=u_0(1-r^2/R^2)$：\fml{Q=\frac{\pi R^2u_0}{2},\;\bar u=\frac{u_0}{2}}

\sect{母题7\;旋转容器}
自由液面：\fml{z=z_0+\frac{\omega^2r^2}{2g},\qquad \Delta z=\frac{\omega^2R^2}{2g}}
无溢出时平均液深守恒：\fml{h_0=z_0+\frac{\omega^2R^2}{4g}}
同一水平面上压差：\fml{p_2-p_1=\rho\frac{\omega^2(r_2^2-r_1^2)}{2}}
\warn{旋转抛物体体积=同高圆柱体积的一半}

\sect{母题8\;伯努利联立连续}
\fml{\frac{p}{\rho g}+\frac{V^2}{2g}+z=\text{const}}
水平变截面：$A_1V_1=A_2V_2$，且\\
\fml{V_2=\sqrt{\frac{2(p_1-p_2)}{\rho[1-(A_2/A_1)^2]}}=\sqrt{\frac{2(p_1-p_2)}{\rho[1-(d_2/d_1)^4]}}}
\fml{Q=A_2V_2}
大水箱小孔：$p_1=p_2=0,\;V_1\approx0$，故\key{托里拆利}：\fml{V=\sqrt{2gH}}

\sect{母题9\;皮托/文丘里/泵功率}
皮托管：\fml{p_0=p+\frac{1}{2}\rho V^2,\qquad V=\sqrt{\frac{2(p_0-p)}{\rho}}}
若U管示液密度$\rho_m$：\fml{V=\sqrt{\frac{2(\rho_m-\rho)g\Delta h}{\rho}}}
泵对流体功率：\fml{N=\rho gQH_p}；若有效率$\eta$，轴功率$N_s=N/\eta$

\sect{母题10\;动量求受力}
\fml{\sum F_x=\rho Q(V_{2x}-V_{1x}),\qquad \sum F_y=\rho Q(V_{2y}-V_{1y})}
平板正冲：\fml{F=\rho QV}；若只改变方向不变速，按速度\key{矢量差}算\\
\warn{先求"壁面对流体"的力，再反号得到"流体对壁/板"的力}

\sect{动量方程符号手术刀（压轴20分）}
\warn{双重正负号法则}——等号左边每项$=$\;[速度方向]\;$\times$\;[质量流量]

\key{法则1}（括号内$\cdot$质量流量）：\warn{只看进出，不看坐标轴！}\\
进水$\to(-\rho uS)$（速度与外法向反向带$-$）；出水$\to(+\rho uS)$（同向带$+$）

\key{法则2}（括号外$\cdot$速度投影）：\warn{严格服从坐标系！}\\
速度投影与正方向同向$\to$正号$+u$；反向$\to$负号$(-u)$

\key{法则3}（等号右边$\cdot$压力）：压力永远\key{推向CV内部}\\
投影到轴，与正方向同向$\to+pS$；反向$\to-pS$

\key{法则4}（约束力）：先猜方向列式，算出负数就是猜反了

\exm{例：三管水缸（$+x$向左，$+y$向上）}
A管水平向左进；B管$45°$左上进；C管竖直向下出

$x$方向动量项：\\
$\overbrace{u_A}^{左=+}\!\!\overbrace{(-\rho u_AS_A)}^{\text{进}=-}\!+\!\overbrace{u_B\cos45°}^{左=+}\!\!\overbrace{(-\rho u_BS_B)}^{\text{进}=-}$\\[0.5pt]
$=\overbrace{P_AS_A}^{\text{推左}=+}\!+\!\overbrace{P_BS_B\cos45°}^{\text{推左}=+}\!-F_2$

$y$方向动量项：\\
$\overbrace{u_B\sin45°}^{上=+}\!\!\overbrace{(-\rho u_BS_B)}^{\text{进}=-}\!+\!\overbrace{(-u_C)}^{下=-}\!\!\overbrace{(\rho u_CS_C)}^{\text{出}=+}$\\[0.5pt]
$=\overbrace{P_BS_B\sin45°}^{\text{推上}=+}\!-F_1$

$F_2$猜向右（$-x$），$F_1$猜向下（$-y$）；正数$=$猜对方向\\
\warn{口诀：里面管进出、外面管方向，互不干涉！}

\sect{母题11\;应力张量/微分方程}
\fml{\vec p_n=\vec n\cdot \mathbf P,\qquad \mathbf P=
\left[\begin{smallmatrix}
p_{xx}&p_{xy}&p_{xz}\\
p_{yx}&p_{yy}&p_{yz}\\
p_{zx}&p_{zy}&p_{zz}
\end{smallmatrix}\right]}
$\mathbf P$对称，仅6个独立分量；静止/理想流体：\fml{\mathbf P=-p\mathbf I}
欧拉：\fml{\rho \frac{D\vec v}{Dt}=-\nabla p+\rho \vec f}
N-S：\fml{\rho \frac{D\vec v}{Dt}=-\nabla p+\mu\nabla^2\vec v+\rho \vec f}
边界：粘性流\key{无滑移}；无粘流只需\key{无穿透}
}

\vspace{0.8pt}\hrule height 0.4pt\vspace{0.8pt}

% ============================================================
\chap{真题套路\;一眼识别}
% ============================================================

\sect{看到题目先想什么}
\key{给速度场}→加速度/流线/有旋/连续性\\
\key{给液面和曲面}→平曲面总压力，先拆$F_x,F_y$\\
\key{给管径压差高度差}→连续+伯努利\\
\key{给进出口方向}→动量方程分$x,y$\\
\key{给旋转角速度}→抛物面自由液面和同高压差\\
\key{给多个液柱高度}→U管爬楼梯法

\sect{最容易丢分的5处}
(1)$h_C$写成斜距而不是垂深；(2)面积/体积忘乘宽度$L$；(3)伯努利跨流线乱用；\\
(4)动量方程没反号；(5)非定常流线没先代$t=t_0$

\sect{三套万能联立}
\key{管流测速}：连续+$p/\rho g+V^2/2g+z$\\
\key{板/弯管受力}：先伯努利求$p$，再动量求$F$\\
\key{液面升降}：RTT质量方程拆\key{储存项+通量项}

\sect{算到最后必检查}
$p$统一用表压还是绝压？$Q$用m$^3$/s？直径mm是否化成m？\\
题目问的是\key{壁对流体}还是\key{流体对壁}？问总压力还是分力？写没写方向？

\sect{定义原句\;考试就这么写}
\key{理想流体}：没有粘性的流体；当粘性力远小于其他力时可近似看作理想流体\\
\key{定常流}：流动参数不随时间变；\key{非定常流}：流动参数随时间变\\
\key{流线}：某瞬时处处与速度矢量相切的曲线；特点：不相交、处处有、定常中与迹线重合\\
\key{压力中心}：总压力作用点；\key{压力体}：受压面外轮廓向上引垂线至液面围成的体积\\
\key{连续介质}：把流体视为连续分布于空间的介质，便于使用微积分

\sect{速度场四连问\;综合题}
若$u=3x^2,\;v=-6xy$，则\\
连续性：$u_x+v_y=6x-6x=0\;\checkmark$\\
加速度：$a_x=uu_x+vu_y=18x^3$，$a_y=uv_x+vv_y=18x^2y$\\
在$(2,1)$：\key{$a_x=144,\;a_y=72$}\\
有旋：\key{$\omega_z=(v_x-u_y)/2=(-6y-0)/2=-3$}\\
变形率：\key{$\varepsilon_{xx}=6x,\;\varepsilon_{yy}=-6x,\;\varepsilon_{xy}=-3y$}

\sect{第四章高频联立\;连续+液面+测速}
\key{大容器液面升降}：\fml{\frac{dh}{dt}=\frac{Q_{in}-Q_{out}}{A_0}
=\frac{A_1V_1-A_2V_2}{A_0}}
若题目给\key{空气不可忽略}，直接列质量式：\\
\fml{(\rho-\rho_a)A_0\frac{dh}{dt}+\rho A_2V_2-\rho A_1V_1=0}
\warn{先判断“升/降”再定正负；若只问瞬时变化率，全部代题给时刻数据}

\sect{测速四件套\;考场一眼套}
\key{托里拆利}：\fml{V=\sqrt{2gH}}
\key{皮托管}：\fml{p_0=p+\frac12\rho V^2,\qquad
V=\sqrt{\frac{2(p_0-p)}{\rho}}}
\key{U管示液测速}：\fml{V=\sqrt{\frac{2(\rho_m-\rho)g\Delta h}{\rho}}}
\key{文丘里/收缩管}：\fml{V_2=\sqrt{\frac{2(p_1-p_2)}
{\rho[1-(d_2/d_1)^4]}} ,\qquad Q=\frac{\pi d_2^2}{4}V_2}
\warn{凡是“管径变化+压差/液柱差”几乎都先联立连续方程}

\sect{动量方程模板\;别再丢符号}
\fml{\sum F_x=\rho Q(V_{2x}-V_{1x}),\qquad
\sum F_y=\rho Q(V_{2y}-V_{1y})}
\key{做题顺序}：先画CV和坐标轴；再把每个进出口速度投影到$x,y$；再把压力力写成\key{推向CV内部}；最后算出的是\key{壁面对流体}，若题目问\key{流体对壁/支座}要反号。\\
\key{常见秒答}：平板正冲\key{$F=\rho QV$}；90$^\circ$转弯要按\key{矢量差}算；若出口对大气且写表压，出口压力项常为\key{0}

\sect{平面/曲面静压\;配合第四章综合}
\key{平面壁}：\fml{F=\rho g h_C A,\qquad
y_D=y_C+\frac{I_C}{y_CA}}
\key{曲面壁}：\fml{F_x=\rho g h_{Cx}A_x,\qquad F_y=\rho gV}
\key{方向判断}：实压力体$\Rightarrow F_y$向下；虚压力体$\Rightarrow F_y$向上；圆柱/圆弧闸门合力线\key{过圆心}

\sect{RTT三件套\;把总公式背成口决}
\key{总公式}：\fml{\frac{dB_{sys}}{dt}
=\frac{\partial}{\partial t}\int_{CV}\beta\rho\,dV
+\int_{CS}\beta\rho\vec V\cdot d\vec A}
\key{质量}：取$B=m,\beta=1$，得\key{储存项+净流出=系统质量变化率}\\
\key{动量}：取$B=\vec P,\beta=\vec V$，得\key{外力=控制体内动量变化+穿过控制面的净动量流}

\sect{平均速度积分\;速度分布题别慌}
\key{定义}：\fml{\bar u=\frac{1}{A}\int_A u\,dA,\qquad Q=\int_A u\,dA}
\key{圆管轴对称}：\fml{Q=\int_0^R u(r)\,2\pi r\,dr}
\key{矩形截面}：\fml{Q=\int_0^b u(y)\,L\,dy}
\key{线性分布}若$u=ky$，则\key{$\bar u=u_{max}/2$}；抛物线分布若$u=u_0(1-r^2/R^2)$，则\key{$\bar u=u_0/2$}

\sect{第四章判断秒杀\;一看就定}
``$\nabla\!\cdot\!\vec v=0$就是定常流''$\Rightarrow$\key{错}，它只表示\key{不可压}\\
``定常流加速度一定为0''$\Rightarrow$\key{错}，迁移加速度仍可能存在\\
``连续方程只适合理想流体''$\Rightarrow$\key{错}，它来自\key{质量守恒}\\
``伯努利方程一定能跨流线''$\Rightarrow$\key{错}，一般仅沿流线；无旋才可跨\\
``动量方程算出的力就是支座受力''$\Rightarrow$\key{不一定}，先分清是\key{壁对流体}还是\key{流体对壁}

\sect{积分结果速翻\;不用现场积分}
\key{线性分布} $u=ky,\;0\le y\le h$：\key{$Q=Lkh^2/2,\;\bar u=kh/2=u_{max}/2$}\\
\key{抛物线分布} $u=u_0(1-r^2/R^2)$：\key{$Q=\pi R^2u_0/2,\;\bar u=u_0/2$}\\
\key{若给质量流量}：\fml{\dot m=\rho Q=\rho A\bar u}
\key{若给重度}：\fml{\gamma=\rho g,\qquad p=\gamma h}

\sect{第四章三行模板\;不会就先写}
\key{质量守恒}：\fml{\frac{\partial}{\partial t}\int_{CV}\rho\,dV+\int_{CS}\rho\vec V\cdot d\vec A=0}
\key{稳流单进单出}：\fml{\dot m_1=\dot m_2}
\key{不可压}：\fml{Q_1=Q_2,\qquad A_1V_1=A_2V_2}

\sect{连续/伯努利10个秒答\;翻到就抄}
(1)流量永远是\key{平均速度乘面积}，不是最大速度乘面积\\
(2)管径越小速度越大；若直径减半且$Q$不变，速度变\key{4倍}\\
(3)大液面常取\key{$V\approx0$}；自由出口表压常取\key{$p=0$}\\
(4)密闭液面常先加\key{$p_0/(\rho g)$}等效液柱\\
(5)皮托管是\key{总压-静压}；U管测速要先写\key{液柱压差}\\
(6)稳流不代表无加速度；仍可能有\key{迁移加速度}\\
(7)不可压二维流反求另一分量时先用\key{$u_x+v_y=0$}\\
(8)已知速度分布求$Q$时先写\key{$Q=\int_A u\,dA$}\\
(9)已知$Q$求受力时先把\key{$V=Q/A$}求出来再动量\\
(10)题目凡是“联立”二字，八成就是\key{连续+伯努利}或\key{伯努利+动量}

\sect{已知量互转\;小块但常用}
\key{已知重度} $\gamma$：先换\key{$\rho=\gamma/g$}\\
\key{已知运动粘度} $\nu$：先换\key{$\mu=\rho\nu$}\\
\key{已知质量流量} $\dot m$：先换\key{$Q=\dot m/\rho$}\\
\key{已知压差水头} $h$：先换\key{$\Delta p=\rho gh$}

\sect{基础碎点补录\;把应急打印并进来}
{\fontsize{4.9}{5.35}\selectfont
\key{连续介质假设}：把流体看作连续分布介质，密度、压强、速度可写成时空连续函数\\
\key{理想流体}：\key{无粘性}；\key{实际流体}：有粘性、可承受切应力\\
\key{牛顿内摩擦定律}：\fml{\tau=\mu \frac{du}{dy}}，切应力与速度梯度成正比\\
\key{绝对压强/表压/真空度}：\fml{p_{abs}=p_a+p_g,\qquad p_{vac}=p_a-p_{abs}}\\
\key{等压面性质}：静止同种连通液体同一水平面压强相等；平衡液面与质量力\key{垂直}\\
\key{欧拉法}看固定空间点；\key{拉格朗日法}跟踪流体质点\\
\key{定常/非定常}看参数是否随时间变；\key{均匀/非均匀}看同一时刻沿程是否变化\\
\key{元流/总流/流束}：流束侧壁上速度与壁面相切，流体\key{不能穿过流束侧壁}
}

\sect{定义原句再补8条\;继续压空白}
{\fontsize{4.75}{5.2}\selectfont
(1)流体质点是\key{宏观上足够小、微观上又包含大量分子}的流体微团\\
(2)不可压缩流动指流体质点在运动中\key{密度近似不变}\\
(3)有旋流动表示流体微团具有\key{刚体式角速度分量}\\
(4)均匀流指同一时刻流场参数沿流程\key{不随空间位置变化}\\
(5)急变流中速度和压强沿程变化\key{很快}，常不能按均匀流处理\\
(6)自由液面上压强通常与外界相通，故表压常取\key{0}\\
(7)压力体是为求曲面壁\key{竖向分力}而引入的辅助体积\\
(8)写题时先判定属于\key{静力、运动学、守恒方程、测速受力}哪一类
}

\sect{名词对照速记\;再补6句}
{\fontsize{4.68}{5.1}\selectfont
(1)表压$=$绝压$-$大气压；\key{真空度}$=$大气压$-$绝压\\
(2)动力粘度$\mu$单位Pa$\cdot$s；运动粘度$\nu$单位\key{m$^2$/s}\\
(3)压力中心是\key{合力作用点}；形心是\key{几何中心}\\
(4)总压$=$静压$+$动压；皮托管前端测\key{总压}\\
(5)控制体关注\key{空间区域}；系统关注\key{同一批流体}\\
(6)微分式偏重\key{局部点}；积分式偏重\key{整个控制体}
}

\sect{超短口决\;第三页补满}
{\fontsize{4.62}{5.0}\selectfont
(1)先分题型，再套方程\\
(2)先画图，再列式\\
(3)先统一单位，再代数值\\
(4)先看表压绝压，再看大气口\\
(5)先判方向，再写正负\\
(6)先求平均速度，再谈流量
}

\vspace{0.8pt}\hrule height 0.4pt\vspace{0.8pt}

\end{multicols}

\clearpage

\begin{multicols}{3}

% ============================================================
\chap{必刷题库\;8题带答案}
% ============================================================

{\fontsize{5.1}{5.7}\selectfont
\exm{1. 流线题}
$u=2x,v=-2y$，求过$(2,2)$的流线。答：$\frac{dx}{2x}=\frac{dy}{-2y}\Rightarrow xy=C$，故\key{$xy=4$}

\exm{2. 加速度题}
$u=2x+t,v=-2y+3$，求$(1,2,1)$处加速度。答：$u=3,v=-1$，\\
\key{$a_x=1+2u=7,\quad a_y=-2v=2,\quad |a|=\sqrt{53}$}

\exm{3. 连续性反求分量}
不可压二维流，$u=3x^2y$且$v|_{y=0}=0$。答：$u_x=6xy$，$v_y=-6xy$，\\
\key{$v=-3xy^2$}

\exm{4. 抛物线速度分布}
$u=u_0(1-r^2/R^2)$，求流量与平均速度。答：\key{$Q=\pi R^2u_0/2,\;\bar u=u_0/2$}

\exm{5. 液面变化率}
容器截面$A_0$，进出流量$Q_1,Q_2$。答：\key{$dh/dt=(Q_1-Q_2)/A_0$}；若$Q_2>Q_1$则液面下降

\exm{6. 侧孔出流}
开口大水箱液面到孔口高差$H$。答：$p_1=p_2=0,V_1\approx0$，\\
\key{$V=\sqrt{2gH}$}

\exm{7. 皮托+U管}
示液密度$\rho_m$，被测流体密度$\rho$，液面差$\Delta h$。答：\\
\key{$V=\sqrt{2(\rho_m-\rho)g\Delta h/\rho}$}

\exm{8. 变截面管流量}
水平管$d_1\to d_2$，压差$\Delta p=p_1-p_2$。答：\\
\key{$V_2=\sqrt{\frac{2\Delta p}{\rho[1-(d_2/d_1)^4]}}$}，\key{$Q=\pi d_2^2V_2/4$}
}

\vspace{0.8pt}\hrule height 0.4pt\vspace{0.8pt}

\chap{填空判断\;考前30分}

{\fontsize{4.9}{5.6}\selectfont
\sect{定义快背}
\key{流体质点}：宏观小、微观大、含大量分子的流体微团\\
\key{连续介质}：把流体视为连续分布介质，物理量可写成空间时间连续函数\\
\key{压力中心}：总压力作用点；\key{压力体}：受压面外轮廓向上引垂线至液面围成体积\\
\key{控制体}：流场中某个确定空间区域；\key{系统}：同一批流体质点

\sect{填空原句}
粘度在普通压强下主要随\key{温度}变；迹线用\key{拉格朗日法}；PIV属\key{欧拉法}\\
流线处处与速度矢量\key{相切}；定常流三线\key{合一}\\
RTT文字：系统量随体导数=控制体内偏$t$+\key{通过控制面的净通量}\\
伯努利7条件：\key{理重定热轴压线}

\sect{判断题秒杀}
``理想流体不可压''→$\times$，理想只表示\key{无粘}\\
``定常流加速度必为0''→$\times$，仍可有迁移加速度\\
``同种液体同深度压强相等''→$\times$，还要\key{连通+静止}\\
``流线可以相交''→$\times$；``圆柱曲面合力线过圆心''→$\checkmark$\\
``$\nabla\!\cdot\!\vec v=0$就是理想流体定常式''→$\times$，那是\key{不可压}条件\\
``伯努利可随便跨流线''→$\times$，除非\key{无旋}

\sect{每章最该背的式子}
\key{Ch1} $\tau=\mu du/dy$，$\nu=\mu/\rho$，水滴$\Delta p=2\sigma/R$，气泡$\Delta p=4\sigma/R$\\
\key{Ch2} $p=p_0+\rho gh$，$F=\rho gh_CA$，$y_D=y_C+I_C/(y_CA)$，$F_y=\rho gV_{\text{压}}$\\
\key{Ch3} $D/Dt=\partial/\partial t+u\partial_x+v\partial_y+w\partial_z$，$\omega_z=(v_x-u_y)/2$\\
\key{Ch4} $\partial_t\rho+\nabla\!\cdot(\rho\vec v)=0$，$A_1V_1=A_2V_2$，$\vec p_n=\vec n\cdot\mathbf P$\\
\key{Ch5} $p/\rho g+V^2/2g+z=\text{const}$，$p_0=p+\rho V^2/2$，$\sum\vec F=\rho Q(\beta_2\vec V_2-\beta_1\vec V_1)$

\sect{大题模板压缩版}
\key{静力学}：先画投影和压力体，再算$F_x,F_y$，最后\key{写方向}\\
\key{伯努利}：先连续后伯努利，优先选\key{液面/出口/已知量多}截面\\
\key{动量}：先求进出口速度和压强，再列$x,y$分量，最后\key{反号}\\
\key{速度场}：先求偏导，再代公式，最后代点值；非定常流线先代$t=t_0$

\sect{临场自测}
能否闭眼写出：$\tau=\mu du/dy$、$p=p_0+\rho gh$、$a_x$公式、$dx/u=dy/v$、RTT、连续式、\\
$F_x/F_y$、旋转液面、伯努利、皮托测速、动量方程、$\vec p_n=\vec n\cdot\mathbf P$？\\
若这12条都能默写，\key{填空判断基本稳}；若8个母题都会算，\key{大题就有90+底子}

\sect{常用数据}
$\rho_w=1000$，$\rho_{Hg}=13600$，$\gamma_w=9810$N/m$^3$，$p_{atm}=101.3$kPa$=10.33$mH$_2$O\\
层流：$\alpha=2,\beta=4/3$；湍流近似：$\alpha\approx1.05,\beta\approx1.02$

\vspace{0.8pt}\hrule height 0.4pt\vspace{0.8pt}

\chap{真题地图\;06/07/08}

\sect{2006真题}
流线+迹线；液面变化率；压力表读数；拖板拉力；旋转容器；复式U管；抛物线速度求流量；\\
曲面闸门总压力；变截面收缩管流量。\\
\key{结论}：06年几乎覆盖了\key{Ch1-5全部硬题型}

\sect{2007真题}
流线过点；三维加速度+有旋；旋转圆柱容器；油管U管+\key{伯努利联立}；泵功率。\\
\key{结论}：07年最爱把\key{连续+伯努利+测压计}揉在一起

\sect{2008真题}
轴承功率；封闭容器压强；带顶盖旋转容器；半球面堵头；\\
连续/加速度/旋转角速度/变形率四连问；流线+迹线；平板截水柱受力。\\
\key{结论}：08年偏爱\key{综合题和细节判定}

\vspace{0.8pt}\hrule height 0.4pt\vspace{0.8pt}

\chap{再刷8结论\;最后冲刺}

\exm{9. 平面总压力中心}
矩形板高$h$、宽$b$、顶边淹深$H$。答：$h_C=H+h/2$，$F=\rho gh_Cbh$，\\
\key{$y_D=y_C+\frac{bh^3/12}{y_C\cdot bh}$}

\exm{10. 曲面总压力}
先算$F_x$和$F_y$，最后$F=\sqrt{F_x^2+F_y^2}$；若为圆弧门，\key{作用线必过圆心}

\exm{11. 封闭容器题}
液面上方气体表压$p_0$不为零时，先把液面抬高\key{$\Delta h=p_0/(\rho g)$}，再做静力学

\exm{12. 旋转容器无溢出}
边缘刚到容器口：先用$z=z_0+\omega^2r^2/(2g)$，再用\key{平均液深不变}解$\omega$

\exm{13. 层流圆管}
$u=u_{max}(1-r^2/R^2)$，平均速度永远\key{$u_{max}/2$}，动能修正\key{$\alpha=2$}

\exm{14. 弯管支反力}
若已知$Q,d,p,\theta$，先算$V$、再算截面压强力、再列动量；\\
求得的是\key{壁对流体}，题目要支座受力就\key{整体反号}

\exm{15. 泵功率}
忽略损失时，泵扬程$H_p$就是两截面总水头差；功率\key{$N=\rho gQH_p$}

\exm{16. 考前最后5分钟}
能否默写：$F=\rho gh_CA$、$F_y=\rho gV$、$dx/u=dy/v$、$a_x$、RTT、连续式、伯努利、\\
皮托测速、动量方程、$\vec p_n=\vec n\cdot\mathbf P$？不能就先背这些

\vspace{0.8pt}\hrule height 0.4pt\vspace{0.8pt}

\chap{口袋问答\;闭眼会写}

\sect{为什么题}
为什么常用欧拉法？答：\key{便于描述整个流场空间分布}。\\
为什么定常流加速度可不为0？答：因\key{迁移加速度}仍可能存在。\\
为什么压力中心在形心下方？答：$y_D=y_C+I_C/(y_CA)$，后项恒正。\\
为什么圆柱曲面合力线过圆心？答：各点压力法线都\key{通过圆心}。\\
为什么大水面速度近似0？答：由连续方程，$A_{tank}\gg A_{hole}$。

\sect{再问5个}
为什么皮托管测总压？答：迎流口\key{驻点}速度降为0。\\
为什么伯努利不能乱跨流线？答：一般只沿流线积分；跨流线需\key{无旋}。\\
为什么$\nabla\!\cdot\!\vec v=0$不是“理想流体条件”？答：它表征\key{不可压}。\\
为什么封闭容器要加假想液面？答：液面上已有气体压强，需先折算成\key{等效液柱}。\\
为什么动量题最后常要反号？答：方程求出多是\key{壁对流体}，题目常问\key{流体对壁}。

\vspace{0.8pt}\hrule height 0.4pt\vspace{0.8pt}

\chap{错题清单\;最后看一遍}

\sect{必查9项}
(1)单位统一：mm$\to$m，kPa$\to$Pa；(2)$p$统一表压/绝压；(3)$g$取9.80还是9.81；\\
(4)$h_C$是不是垂深；(5)面积体积是否乘宽度$L$；(6)流线题是否先代$t=t_0$；\\
(7)连续性是否联立；(8)动量力方向是否标清；(9)答案有没有\key{大小+方向+作用线}

\vspace{0.8pt}\hrule height 0.4pt\vspace{0.8pt}

\chap{12个一眼答案}

\sect{临场秒扫}
直径缩半$\Rightarrow$速度\key{4倍}；层流圆管$\bar u=u_{max}/2$；静止流体切应力\key{为0}\\
自由射流出口表压\key{为0}；壁面切应力\key{最大在壁面}；圆柱曲面合力线\key{过圆心}\\
定常流中流线/迹线/脉线\key{合一}；大水箱液面速度\key{约为0}；皮托管测\key{总压}\\
液体$T\uparrow\Rightarrow \mu\downarrow$；气体$T\uparrow\Rightarrow \mu\uparrow$；粘性流固壁\key{无滑移}
}

\vspace{0.8pt}\hrule height 0.4pt\vspace{0.8pt}

\chap{第四章再压缩\;最后塞满}

{\fontsize{4.75}{5.35}\selectfont
\sect{旋转液面\;会公式就会题}
\fml{z=z_0+\frac{\omega^2r^2}{2g},\qquad
\Delta z=\frac{\omega^2R^2}{2g}}
\fml{h_0=z_0+\frac{\omega^2R^2}{4g},\qquad
p_2-p_1=\rho\frac{\omega^2(r_2^2-r_1^2)}{2}}
\key{核心两句}：自由液面是\key{旋转抛物面}；无溢流时用\key{平均液深不变}，刚好到口时再加边界条件

\sect{伯努利联立模板\;最该背}
\key{标准式}：\fml{\frac{p_1}{\rho g}+\frac{V_1^2}{2g}+z_1
=\frac{p_2}{\rho g}+\frac{V_2^2}{2g}+z_2}
\key{联立顺序}：先连续$A_1V_1=A_2V_2$，再伯努利，最后代$Q=AV$\\
\key{截面优先}：大液面/自由出口/已知压强或高度多的截面优先；大水箱常取\key{$V_1\approx0$}，自由射流出口表压常取\key{$p=0$}

\sect{弯管受力\;一页模板}
\key{第1步} 求截面速度：\fml{V=Q/A}
\key{第2步} 若未知压力，先用伯努利求$p_1,p_2$\\
\key{第3步} 分$x,y$列动量式：\fml{\sum F_x=\rho Q(V_{2x}-V_{1x})}
\fml{\sum F_y=\rho Q(V_{2y}-V_{1y})}
\key{第4步} 支座受力$=$流体对壁$=$\key{壁对流体反号}

\sect{考前必背16条\;闭眼默写}
(1)$p=p_0+\rho gh$\quad
(2)$F=\rho gh_CA$\quad
(3)$y_D=y_C+I_C/(y_CA)$\\
(4)$F_x=\rho gh_{Cx}A_x$\quad
(5)$F_y=\rho gV$\quad
(6)$dx/u=dy/v=dz/w$\\
(7)$dx/dt=u,\;dy/dt=v$\quad
(8)$$\displaystyle \frac{D}{Dt}=\partial_t+u\partial_x+v\partial_y+w\partial_z$$\\
(9)$a_x=u_t+uu_x+vu_y+wu_z$\quad
(10)$\omega_z=(v_x-u_y)/2$\\
(11)$\partial_t\rho+\nabla\!\cdot(\rho\vec v)=0$\quad
(12)$A_1V_1=A_2V_2$\\
(13)$V=\sqrt{2gH}$\quad
(14)$p_0=p+\rho V^2/2$\quad
(15)$\sum F_x=\rho Q(V_{2x}-V_{1x})$\quad
(16)$\vec p_n=\vec n\cdot \mathbf P$

\sect{最后一分钟\;直接翻这里}
\key{看到液柱高差}先想U管爬楼；\key{看到曲面闸门}先拆$F_x,F_y$；\key{看到容器液面变化}先写$dh/dt$；\\
\key{看到管径变化+压差}先连续再伯努利；\key{看到弯管/喷嘴受力}先速度投影再动量；\\
\key{看到旋转容器}先抛物面；\key{看到速度场}先偏导，一口气连问连续/加速度/旋度/变形率

\sect{8个母题再压一遍\;就背这几句}
\key{密闭容器静压}：先把气体表压折成\key{$\Delta h=p_0/(\rho g)$}，再按静水压算\\
\key{矩形平板闸门}：$h_C=H+h/2$，$I_C=bh^3/12$，故\key{$y_D=y_C+h^2/(12y_C)$}\\
\key{圆弧/曲面闸门}：先求\key{$F_x$投影、$F_y$压力体}，再合成\key{$F=\sqrt{F_x^2+F_y^2}$}\\
\key{大水箱小孔出流}：\key{$V=\sqrt{2gH}$}，大液面速度近似0\\
\key{测速U管}：先写\key{$p_0-p=(\rho_m-\rho)g\Delta h$}，再接\key{$p_0-p=\rho V^2/2$}\\
\key{变截面流量}：先\key{$A_1V_1=A_2V_2$}，再伯努利\\
\key{弯管受力}：先速度、再压力、后动量，最后别忘\key{反号}\\
\key{旋转液面}：先抛物面，后平均液深不变，再代边界

\sect{12句判断题原句\;会考口吻}
理想流体\key{不等于}不可压流体；不可压流满足\key{$\nabla\cdot\vec v=0$}\\
流线在同一瞬时处处与速度矢量\key{相切}；定常流中流线/迹线/脉线\key{重合}\\
静止流体\key{不能承受切应力}；静止流体中同一点各方向压强\key{相等}\\
压力中心一般在形心\key{下方}；圆柱曲面合力线\key{过圆心}\\
皮托管迎流端是\key{驻点}；自由射流出口表压通常\key{为0}\\
伯努利方程不是万能式，必须先看\key{条件和选面}

\sect{考场保分模板\;不会也先写}
\key{连续方程题}：写$Q=AV$或$\partial_t\rho+\nabla\!\cdot(\rho\vec v)=0$，再说明\key{对不可压稳流可化简}\\
\key{伯努利题}：先写两截面、三项水头、条件“稳、不可压、沿流线、忽略损失”\\
\key{动量题}：先写控制体、坐标轴、进出口速度投影、压力方向、再列$x,y$分量\\
\key{静压力题}：先找形心深度$h_C$，平面壁写$F=\rho gh_CA$；曲面壁写$F_x,F_y$后合成

\sect{典型数字\;抄上就能省时间}
$\rho_w=1000$ kg/m$^3$，$\rho_{Hg}=13600$ kg/m$^3$，$g=9.81$ m/s$^2$\\
$1$ atm$=101.3$ kPa$=10.33$ mH$_2$O；水的重度$\gamma_w=9810$ N/m$^3$\\
层流圆管：\key{$\alpha=2,\beta=4/3$}；湍流工程近似：\key{$\alpha\approx1.05,\beta\approx1.02$}

\sect{最后保命清单\;落笔前扫一眼}
单位有没有全化成SI；直径半径有没有混；mm有没有换m；kPa有没有换Pa\\
深度$h_C$是不是\key{垂直深度}；面积体积有没有乘宽度；结果有没有写\key{大小+方向+作用点/作用线}\\
题目问的是\key{壁对流体}还是\key{流体对壁}；出口是不是大气；表压和绝压有没有混

\sect{20个秒套例题\;只看结论}
{\fontsize{4.55}{5.05}\selectfont
(1)矩形竖板顶边在液面：\key{$F=\rho gbh^2/2,\;y_D=2h/3$}\\
(2)矩形竖板顶边距液面$H$：\key{$F=\rho gbh(H+h/2)$}\\
(3)密闭容器先折液柱：\key{$h_e=p_0/(\rho g)$}\\
(4)大水箱小孔出流：\key{$V=\sqrt{2gH}$}\\
(5)密闭大水箱小孔：\key{$V=\sqrt{2(p_0/\rho+gH)}$}\\
(6)皮托测速：\key{$V=\sqrt{2\Delta p/\rho}$}\\
(7)U管测速：\key{$V=\sqrt{2(\rho_m-\rho)g\Delta h/\rho}$}\\
(8)收缩管流量：\key{$Q=A_2\sqrt{\frac{2(p_1-p_2)}{\rho[1-(A_2/A_1)^2]}}$}\\
(9)液面升降：\key{$dh/dt=(Q_1-Q_2)/A_0$}\\
(10)直径减半且$Q$不变：\key{$V$变4倍}\\
(11)射流正冲平板：\key{$F=\rho QV$}\\
(12)射流完全反向：\key{$F=2\rho QV$}\\
(13)抛物线流速分布：\key{$Q=\pi R^2u_0/2,\;\bar u=u_0/2$}\\
(14)不可压二维反求$v$：由\key{$u_x+v_y=0$}积分\\
(15)旋转液面高差：\key{$\Delta z=\omega^2R^2/(2g)$}\\
(16)速度场加速度：先求各偏导再代\key{$a_x,a_y$}\\
(17)流线题：先定时刻，再写\key{$dx/u=dy/v=dz/w$}\\
(18)迹线题：直接写\key{$dx/dt=u,\;dy/dt=v$}\\
(19)理想流体应力张量：\key{$\mathbf P=-p\mathbf I$}\\
(20)圆柱/圆弧曲面：合力作用线\key{过圆心}
}

\sect{考前填空20句\;能捡一堆分}
{\fontsize{4.45}{4.95}\selectfont
(1)流线处处与速度矢量\key{相切}\quad
(2)定常流中流线/迹线/脉线\key{重合}\\
(3)静止流体切应力\key{为0}\quad
(4)同一种静止连通液体同一水平面压强\key{相等}\\
(5)压力中心通常在形心\key{下方}\quad
(6)圆柱曲面合力作用线\key{过圆心}\\
(7)大水箱液面速度通常近似\key{0}\quad
(8)自由射流出口表压通常为\key{0}\\
(9)皮托管测的是\key{总压}\quad
(10)驻点速度等于\key{0}\\
(11)不可压条件是\key{$\nabla\cdot\vec v=0$}\quad
(12)连续方程来源于\key{质量守恒}\\
(13)RTT把\key{系统}和\key{控制体}联系起来\\
(14)伯努利方程本质是\key{机械能守恒}\\
(15)动量方程本质是\key{牛顿第二定律}\\
(16)理想流体应力张量\key{$=-p\mathbf I$}\\
(17)粘性流体贴壁满足\key{无滑移}\quad
(18)液体温度升高粘度通常\key{减小}\\
(19)气体温度升高粘度通常\key{增大}\quad
(20)单位一定统一到\key{SI}
}

\sect{应急凑分写法\;最后手段}
{\fontsize{4.45}{4.95}\selectfont
问静压力就先写\key{$p=p_0+\rho gh$}、\key{$F=\rho gh_CA$}；问曲面就补\key{$F_x,F_y$}\\
问流量就先写\key{$Q=AV$}；问测速就先写\key{$p_0-p=\rho V^2/2$}\\
问液面变化就先写\key{$dh/dt=(Q_{in}-Q_{out})/A_0$}\\
问受力就先画控制体、写\key{$\sum F_x,\sum F_y$}；问速度场就先求\key{偏导}
}

\sect{最后再背8句\;压轴收尾}
{\fontsize{4.45}{4.95}\selectfont
(1)连续方程管的是\key{质量守恒}\quad
(2)伯努利方程管的是\key{机械能守恒}\\
(3)动量方程管的是\key{受力与速度变化}\quad
(4)RTT负责把\key{系统}换成\key{控制体}\\
(5)压力体只为求\key{$F_y$}，不是额外真实液体\\
(6)压力中心是\key{总压力作用点}，不等于形心\\
(7)静止流体不能承受\key{切应力}\quad
(8)粘性流贴壁满足\key{无滑移}
}

\sect{落笔前再扫10条\;真最后}
{\fontsize{4.35}{4.8}\selectfont
(1)$Q=AV$里$V$默认\key{平均速度}\quad
(2)$h_C$默认\key{垂直深度}\\
(3)自由出口表压多取\key{0}\quad
(4)大液面速度多取\key{0}\\
(5)求支座受力常要把\key{壁对流体}反号\\
(6)动量题压力力方向永远\key{推向CV内}\\
(7)流线题先\key{定时刻}\quad
(8)迹线题先\key{列微分方程}\\
(9)封闭容器别忘\key{等效液柱}\quad
(10)所有答案尽量写\key{大小+方向}
}

\sect{基础定义补充\;这张图里的高频句}
{\fontsize{4.25}{4.7}\selectfont
(1)液体通常近似看作\key{不可压缩}；气体压缩性\key{明显}\\
(2)液体$T\uparrow$通常$\mu\downarrow$；气体$T\uparrow$通常$\mu\uparrow$\\
(3)静止流体\key{不能承受拉力和切力}，只能承受压应力\\
(4)压力总是沿\key{受压面内法线}方向指向受压面\\
(5)同一点压强在各方向\key{相等}是流体静压基本性质\\
(6)流线是某瞬时速度场的几何图像；迹线是单个质点运动轨迹\\
(7)定常流中\key{流线=迹线=脉线}；非定常流一般不重合\\
(8)一维流/二维流/三维流按速度分量和坐标变化个数区分\\
(9)渐变流可近似按\key{均匀流}处理；急变流局部变化剧烈\\
(10)层流与紊流常用\key{雷诺数}判别；圆管中$Re<2000$常近似层流，$Re>4000$常近似紊流\\
(11)有压流流满管道；无压流有\key{自由液面}\\
(12)总流能量按整个过流断面讨论；元流可看作\key{无穷小流束}\\
(13)分析题目先分清是\key{定义题、公式题、联立题、受力题}再下手\\
(14)不会做时先把\key{基本定义+控制方程+已知量互转}写出来
}

\sect{判断原句再补12条\;继续填满}
{\fontsize{4.18}{4.62}\selectfont
(1)理想流体\key{可以}是可压缩的，理想只表示\key{无粘性}\\
(2)表压为负不代表压强为负，只表示其\key{低于大气压}\\
(3)真空度越大，绝对压强越\key{小}\\
(4)静止液体内部不存在\key{剪切变形率}\\
(5)等压面与质量力垂直，重力场中静止液体等压面通常与自由液面\key{平行}\\
(6)流线不能相交，否则同一点速度方向将\key{不唯一}\\
(7)总流伯努利式中的速度水头常取\key{断面平均速度}计算\\
(8)动量修正系数$\beta$与速度分布有关，速度越不均匀通常$\beta$越\key{大}\\
(9)动能修正系数$\alpha$在层流中比紊流\key{更大}\\
(10)曲面壁总压力通常要先求\key{水平分力和竖直分力}再合成\\
(11)控制体可以固定不动，也可以\key{运动或变形}\\
(12)考试中若题目没特别说明，水默认取$\rho\approx$\key{1000 kg/m$^3$}
}

\sect{一眼判断再补10条\;压到最后}
{\fontsize{4.1}{4.55}\selectfont
(1)等压面\key{不是}一定水平，只有重力单独作用且液体静止时才近似水平\\
(2)速度分布越不均匀，平均速度与最大速度差得越\key{远}\\
(3)层流圆管中最大速度是平均速度的\key{2倍}\\
(4)若闸门两侧都是液体，要算的是\key{两侧压力差的合力}\\
(5)封闭容器静力题往往先做\key{压强折算}再做静水压\\
(6)喷嘴/弯管受力题最怕漏掉\key{压力项和反号}\\
(7)曲面壁竖向分力大小等于压力体内液体\key{重量}\\
(8)控制体法最适合处理\key{进出口流动}问题\\
(9)速度场题先求偏导，往往能一口气做出\key{连续、加速度、旋度、变形率}\\
(10)不会时先写\key{定义+基本方程+代入已知}，比空着强
}

\sect{机械61补丁\;你说不熟的这几项}
{\fontsize{4.05}{4.45}\selectfont
\key{质点导数}：\fml{\frac{DB}{Dt}=\frac{\partial B}{\partial t}
+u\frac{\partial B}{\partial x}+v\frac{\partial B}{\partial y}+w\frac{\partial B}{\partial z}}
\key{定常流}时流线与迹线重合；一维稳定流也常可直接看作\key{重合}\\
\key{速度势函数}：无旋流可设\key{$u=\varphi_x,\;v=\varphi_y,\;w=\varphi_z$}\\
\key{流函数}：二维不可压流可设\key{$u=\psi_y,\;v=-\psi_x$}\\
\key{相对平衡}：自由液面总与\key{合加速度}垂直；水平加速容器常用\key{$\tan\theta=a/g$}\\
\key{连续介质条件}：\key{$Kn=\lambda/L\ll1$}\\
\key{雷诺输运公式}：系统量变化率$=$控制体内积累率$+$穿过控制面的净输出\\
\key{Mach速记}：$Ma<0.3$低速；$0.3\sim0.8$亚音速；$0.8\sim1.2$跨音速；$1.2\sim5$超音速；$>5$高超音速
}
}

\sect{最后口决\;第四页收边}
{\fontsize{4.0}{4.4}\selectfont
(1)伯努利先选面，动量先定轴\\
(2)曲面先拆$F_x,F_y$，平面先找$h_C$\\
(3)大水箱取$V\approx0$，自由出口取$p\approx0$\\
(4)问支座受力，多半最后要反号\\
(5)速度场先偏导，静力题先画深度\\
(6)不会整题，先把定义和主方程写上
}

\end{multicols}
\end{document}
